\section{Introduction}\label{sec:introduction}

Fiat-backed stablecoins are tokens issued against a reserve of conventional dollar assets, designed to be worth one dollar each at every point in time. By 2026 they routinely settle hundreds of billions of dollars across exchanges, cross-border transfers, and on-chain payments. Their economic appeal as a means of payment rests on the simple promise of par convertibility into the dollar. The promise is sometimes briefly broken. In March 2023, after Circle disclosed that part of its cash reserve was held as uninsured deposits at Silicon Valley Bank, the secondary-market price of USDC fell as low as ninety cents per dollar within hours, holders rushed to convert at the issuer, and the price discount spread to DAI through the latter's USDC collateral. The peg was restored within four days. The episode revealed that the same token can trade within a few hundredths of a cent from par for ninety-nine percent of the hours of a year and then become a state-contingent claim worth ninety cents for the remaining one percent.

This paper asks when a fiat-backed stablecoin behaves as the par-convertible payment instrument it promises to be, when it does not, and what the answer implies for using the stablecoin instead of an established payment service in a given corridor. We build a microstructure model of the user's coordinated choice between converting at the issuer and selling on a secondary venue. The choice is a Morris-Shin coordination game with heterogeneous user types: each type observes a noisy private signal about the issuer's reserve adequacy and chooses to convert when the signal is bad enough. The strategic stage is solved as a fixed point on the type-specific thresholds. We calibrate the model to the March 2023 USDC episode in hourly data, evaluate the calibrated mechanism out of sample on USDT and DAI, falsify it on the May 2022 TerraUSD collapse, and use it to compute a per-corridor cost comparison against an incumbent payment service.

The paper has three contributions. The first is theoretical. We characterise the threshold profile of the heterogeneous-type fixed point and prove three properties of the solution. The fixed point exists and is unique under regularity conditions on the signal noise and the secondary-market block. Tighter issuer reserves reduce aggregate conversion pressure at the fixed point, because the anticipated price discount on the exchange makes individual conversion less attractive at the margin. A higher payment-use value of the stablecoin lowers the equilibrium conversion mass. The two comparative-statics results respectively formalise in the heterogeneous-type setting the structural mechanism that \citet{ma_zeng_zhang_2025} document empirically across stablecoins, and the run-attenuation channel that \citet{bertsch_2025_adoption_fragility} identifies at the adoption stage. Proofs are in Appendix~\ref{sec:proofs}.

The second contribution is empirical. We calibrate the model to six moments of the March 2023 USDC episode using hourly data and the regime-dependent dispersion of the latent fundamental as the structural device that separates tranquil and stressed states. The fit is tight on the average severity of the stress regime and on the frequency of large secondary-market dislocations; it is loose on the deepest minute of the realised episode. We then assess the calibrated structural parameters out of sample on USDT and DAI. The parameters transport cleanly to DAI, consistent with DAI's exposure to USDC through its collateral. They do not transport to USDT, whose stable trajectory in the March 2023 window is not reproducible at the calibrated parameters. We read the mismatch as a boundary statement on the model's structural domain: the framework describes stablecoins with open primary-market arbitrage and reserve-backed conversion, and does not describe stablecoins with concentrated arbitrage access. The boundary aligns with the structural difference between USDC and USDT that \citet{ma_zeng_zhang_2025} document empirically. We falsify the model on TerraUSD, which it cannot reproduce because TerraUSD lacked the reserve mechanism on which the fixed point depends.

The third contribution is to apply the calibrated model to a per-corridor cost comparison against an incumbent payment service. For each of three named corridors (US to Mexico, US to the Philippines, Europe to Brazil) we compute a per-unit certainty-equivalent cost for the stablecoin and for the corresponding incumbent service using fee data from the World Bank Remittance Prices Worldwide quarterly panel. We use the metric to evaluate the November 2025 Banco Central do Brasil package that places virtual assets within the foreign-exchange framework, and we report the implicit regulatory burden that would equate the two costs in the cross-border corridor. The metric is a static input to a dynamic analysis of adoption across payment services, which we describe but do not estimate.

We read fiat-backed stablecoins as payment instruments in a segmented rather than a universal way. The cost comparison turns negative for the stablecoin in retail and cross-border corridors, by a margin large enough to survive the calibrated tail risk and a tenfold increase in its empirical frequency. The same comparison turns positive in the wholesale corridor, conditional on a settlement-finality assumption that is reasonable for users who require atomic execution. The segmentation reflects a class-specific congestion cost that captures the institutional fact that wholesale users do not tolerate uncertain settlement, rather than a class-specific run risk that would be identified inside the model. The calibrated model does not deliver large differences in the per-class conversion rate, and we are explicit about this in the body.

Section~\ref{sec:literature} positions the paper. Section~\ref{sec:model} develops the model. Section~\ref{sec:data_calibration} describes the data and the calibration strategy. Section~\ref{sec:results} reports the calibrated parameter vector, the moment-matching diagnostics, the out-of-sample evaluation on USDT and DAI, and the falsification on TerraUSD. Section~\ref{sec:cost_benefit} computes the per-corridor cost comparison and its response to the Brazilian package. Section~\ref{sec:discussion} discusses the implications for payment use cases and points to the dynamic adoption extension. Section~\ref{sec:conclusion} concludes. Appendix~\ref{sec:proofs} contains the proofs of the three propositions.
